% K=1 on the SAME six sites as the 3IA poster.
% Source: Présentations/Posters/3IA_Days_2026/test_six_point_geometry.py.
% Coordinates use r0=1.2 and h=r0*sqrt(3).
% Every shortest edge has length 2*r0=2.4. A and B connect through C;
% D,E,F form the second triangle; C--D joins them at the SAME radius.
% Thus all six components merge in ONE closed plateau at r0=1.2.
% Required: tikz, xcolor; colours inriablue, inriared, inriagrayblue.
\begin{tikzpicture}[line cap=round,line join=round,
  every node/.style={font=\small},
  point/.style={circle,fill=inriagrayblue,inner sep=2.2pt}]
\useasboundingbox (-3.1,-4.0) rectangle (10.5,2.7);
\pgfmathsetmacro{\hx}{1.2*sqrt(3)}
\foreach \shift/\rad/\colour/\caption in {
  0/1.1/inriablue/{Six groupes séparés},
  7.3/1.25/inriared/{Un seul groupe}}
{
  \begin{scope}[xshift=\shift cm]
    \node[font=\small\bfseries,text=\colour] at (0,2.35) {\caption};
    \ifdim\shift pt=0pt
      \node[font=\small,text=inriagrayblue] at (0,1.99) {$r=1{,}1$};
    \else
      \node[font=\small,text=inriagrayblue] at (0,1.99) {$r=1{,}25$};
    \fi
    \begin{scope}[scale=.57]
      \coordinate (A) at (-1.2-\hx,1.2);
      \coordinate (B) at (-1.2-\hx,-1.2);
      \coordinate (C) at (-1.2,0);
      \coordinate (D) at (1.2,0);
      \coordinate (E) at (1.2+\hx,1.2);
      \coordinate (F) at (1.2+\hx,-1.2);
      % Filled UNION; opaque colour keeps overlaps at the same shade.
      \foreach \p in {A,B,C,D,E,F} {\fill[\colour!17] (\p) circle[radius=\rad];}
      \foreach \p in {A,B,C,D,E,F} {\draw[\colour!65,line width=.65pt] (\p) circle[radius=\rad];}
      \foreach \p in {A,B,C,D,E,F} {\node[point] at (\p) {};}
      \foreach \p in {A,E} {\node[above=3pt,text=inriagrayblue] at (\p) {\p};}
      \foreach \p in {B,F} {\node[below=3pt,text=inriagrayblue] at (\p) {\p};}
      \node[above=3pt,text=inriagrayblue] at (C) {C};
      \node[above=3pt,text=inriagrayblue] at (D) {D};
    \end{scope}
  \end{scope}
}
% Dendrogramme non binaire: une seule multifusion, aucun ordre artificiel.
\node[font=\small\bfseries,text=inriagrayblue] at (3.65,-1.95)
  {Quand les disques se touchent, les groupes se rejoignent.};
\draw[inriared,line width=1.2pt] (1.05,-2.65)--(6.25,-2.65);
\foreach \xx/\label in {1.05/A,2.09/B,3.13/C,4.17/D,5.21/E,6.25/F} {
  \draw[inriablue,line width=1pt] (\xx,-3.55)--(\xx,-2.65);
  \node[text=inriagrayblue,below=2pt] at (\xx,-3.55) {\label};
}
\node[text=inriared,font=\small\bfseries,above=3pt] at (3.65,-2.65) {$r=1{,}2$};
\draw[inriagrayblue,->] (.4,-3.6)--(.4,-2.4);
\node[text=inriagrayblue,left=4pt,font=\footnotesize] at (.4,-3.02) {$r$};
\end{tikzpicture}
